% HIKET -- WORKING DOCUMENT: parameterization assumptions, August 2026.
%
% Purpose: state, in one consistent package and in near-manuscript form, every
% assumption that goes into the HIKET parameterization -- what each parameter IS,
% where its prior comes from, and what constrains it. Written after the 2026-08-07
% diagnostic session, which found that several of these assumptions are mutually
% inconsistent. Sections 2-3 describe what is CURRENTLY implemented; section 4
% states the defects; section 5 proposes the replacement package; section 6 derives
% the one genuinely new number. Nothing here is implemented yet.
\documentclass[11pt]{article}
\usepackage[utf8]{inputenc}\usepackage[T1]{fontenc}
\usepackage[margin=2.4cm]{geometry}
\usepackage{graphicx}\usepackage{amsmath}\usepackage{amssymb}\usepackage{xcolor}
\usepackage{booktabs}\usepackage{longtable}
\usepackage{float}\usepackage[hidelinks]{hyperref}
\usepackage{caption}\captionsetup{font=small,labelfont=bf}
\usepackage[most]{tcolorbox}

\definecolor{decbg}{HTML}{EAF7EE}\definecolor{decrule}{HTML}{2E8B57}
\newtcolorbox{proposal}[1]{colback=decbg, colframe=decrule, boxrule=0pt, leftrule=3pt,
  arc=1pt, left=8pt, right=8pt, top=5pt, bottom=5pt, fontupper=\small,
  title={\bfseries PROPOSAL --- #1}, coltitle=decrule, fonttitle=\small\bfseries,
  attach title to upper=\par, before skip=9pt, after skip=9pt}

\definecolor{cavbg}{HTML}{FDF3E7}\definecolor{cavrule}{HTML}{C87A2F}
\newtcolorbox{caveat}{colback=cavbg, colframe=cavrule, boxrule=0pt, leftrule=3pt,
  arc=1pt, left=8pt, right=8pt, top=5pt, bottom=5pt, fontupper=\small,
  before skip=9pt, after skip=9pt}

\definecolor{defbg}{HTML}{FBEAEA}\definecolor{defrule}{HTML}{B03A3A}
\newtcolorbox{defect}[1]{colback=defbg, colframe=defrule, boxrule=0pt, leftrule=3pt,
  arc=1pt, left=8pt, right=8pt, top=5pt, bottom=5pt, fontupper=\small,
  title={\bfseries DEFECT --- #1}, coltitle=defrule, fonttitle=\small\bfseries,
  attach title to upper=\par, before skip=9pt, after skip=9pt}

\newcommand{\sinit}{\sigma_{\mathrm{init}}}
\newcommand{\sinput}{\sigma_{\mathrm{input}}}
\newcommand{\sobs}{\sigma_{\mathrm{obs}}}
\newcommand{\Jfull}{\bar{J}_{\mathrm{full}}}
\newcommand{\Jto}{\bar{J}_{t_0}}

\title{HIKET --- parameterization assumptions\\
  \large Working document: what each parameter is, where its prior comes from,\\
  and a proposed consistent package}
\author{}
\date{7 August 2026 --- working document, not for circulation}

\begin{document}
\maketitle


% ---------------------------------------------------------------------------
% STATUS BLOCK inserted 2026-08-10. Delete once the run below has landed and
% the affected passages have been rewritten.
% ---------------------------------------------------------------------------
\begin{center}
\fbox{\begin{minipage}{0.93\textwidth}
\small
\textbf{DRAFT STATUS --- 2026-08-10.} \textbf{Defect D4 (the likelihood conflates
observation with model error) is now quantified and is being addressed.} The measured log-residual
spread is 0.708--0.735 across all six models against an assumed
$\sigma_{\mathrm{obs}}=0.442$ --- a factor 1.6 --- so the implied \emph{model} error (0.55--0.59)
exceeds the observation error. Because the level penalty scales as $1/\sigma^{2}$, this amplified
the pressure to match stock \emph{levels} by ${\sim}2.6\times$ relative to the priors, which is a
candidate driver of the short calibrated residence times. Roihu jobs 563524--563529 (commit
\texttt{3b0d533}) test this with a total error scale of 0.72, the plug-in estimate; the fraction
prior is deliberately left unchanged so the error model is tested as one factor.
\end{minipage}}
\end{center}

\begin{abstract}
\noindent This document states every parameterization assumption in the HIKET calibration in
one place, in the form they would take in a Methods section, and proposes a revised package.
It was written after a diagnostic session (2026-08-07) on the first recalibration to include
all round-1 changes and both upstream data corrections. That session found the assumptions to
be individually defensible but \emph{mutually inconsistent}: the same external evidence (the
national growing-stock record) is used to set the \emph{shape} of the pre-run litter ramp and
simultaneously contradicted by the prior that sets its \emph{level}; the two ends of the pre-run
are referenced to different litter aggregates, so the parameter that controls the initial state
does not mean what its name and documentation say; and the constraint introduced to keep litter
fluxes physical bounds a quantity orthogonal to the failure it is now being asked to prevent.
Sections~\ref{sec:principle}--\ref{sec:current} describe what is implemented;
section~\ref{sec:defects} states the defects; section~\ref{sec:package} proposes the
replacement; section~\ref{sec:R} derives the one new quantity it requires. A separate,
unresolved problem --- a uniform $\sim$20\,\% over-prediction shared by all six models --- is
recorded in section~\ref{sec:open} because it conditions how the proposal should be read. A paired pre-flight (plain priors) additionally shows the bounded transform to be removable without cost, allowing the whole parameterization to collapse to four ordinary priors.
\end{abstract}

\begin{caveat}
\textbf{Naming.} Earlier drafts numbered the proposals P1--P5. They are now referred to by name,
because the numbering repeatedly caused confusion with the model names (TP2/TP3) and with the
C1--C5 revision series. The mapping, for cross-reference with older notes: P1 = \textsc{common
anchor}; P2 = \textsc{ratio coordinate} (now superseded); P3 = \textsc{ratio prior};
P4 = \textsc{plain priors}; P5 = \textsc{free slow rate}.
\end{caveat}


% =====================================================================
\section{The principle a prior has to satisfy here}\label{sec:principle}

The HIKET priors were built on a three-tier homogenization scheme whose stated aim is that
priors be homogeneous in \emph{construction method, scale, constraint and degrees of freedom}
--- not in numerical value. Tier~1 (climate, size) takes each model's centre from its own
published calibration; Tier~2 (transfer fractions) applies a common weakly-informative logit
prior; Tier~3 was the two auxiliary litter parameters, treated as nuisance quantities with an
identical weak prior in every model.

That scheme has since been extended piecemeal. The auxiliary parameters acquired a physical
window (the \texttt{flux\_pair} transform, July 2026) and one of them acquired a physical
centre (C4b), without a stated rule governing when a Tier-3 nuisance becomes a physical
quantity. The result is the inconsistency this document addresses.

The rule we propose to adopt explicitly is:

\begin{quote}
\emph{A parameter with a physical referent takes its prior from external evidence about that
referent, expressed in that referent's own units. A parameter without a physical referent
should not be given an informative prior at all --- and if it cannot be given one, that is a
reason to reparameterize until the sampled quantity does have a referent.}
\end{quote}

The second clause is the operative one. Both auxiliary parameters currently carry priors that
are informative in practice, and one of them is a derived ratio whose prior cannot be stated in
physical units at all.

\paragraph{A sharper operational version (added 2026-08-07).} The rule above is about
\emph{referents}. In practice what determines how much a prior matters is the \emph{curvature of
the likelihood along it}, and the two auxiliary parameters sit in opposite regimes:

\begin{table}[H]\centering\small
\begin{tabular}{lrrl}
\toprule
 & likelihood cost of a 15\,\% move & prior cost & decided by\\
\midrule
$\sinput$ & $69.3$ nats & $0.62$ nats & \textbf{likelihood} ($\approx 112:1$)\\
$\sinit$ / $R$ & weak, and concentrated in the 1985 obs & --- & \textbf{jointly} (see \S\ref{sec:cond1985})\\
\bottomrule
\end{tabular}
\end{table}

\noindent So: \emph{a prior needs justification in proportion to how flat the likelihood is along
it.} $\sinput$ can carry a weak or even stale justification --- the data overrules it by two
orders of magnitude, which is why re-centring it from $1.30$ to $1.00$ leaves the posterior at
$1.308$ either way. $R$ needs a defensible one, because the likelihood constrains it far more weakly ---
though \emph{not} negligibly, see \S\ref{sec:cond1985}. This explains why the initialisation
parameters have been the difficult ones throughout, and it means effort spent justifying C4b's
$\sinput$ centre is largely wasted while effort spent on $R$ is not.

% =====================================================================
\section{What is currently implemented}\label{sec:current}

\begin{caveat}
\textbf{Everything in this section describes code that runs today} --- it is the parameterization
behind the \texttt{20260805} calibration, not a proposal. The proposed changes are in
section~\ref{sec:package} and are clearly marked as such. Section~\ref{sec:defects} is the bridge:
it states what is wrong with what follows here.
\end{caveat}

\subsection{Decomposition rates}
Yasso07/15/20 do not fit their decomposition rates: the $a$-vector is fixed at published
values and injected as a constant, and what is calibrated is the twelve lateral transfer
fractions, the climate parameters, woody size, and the two auxiliary $\sigma$'s.

C1 (round~1) made the simple models homogeneous with that treatment. The fast rate is
\textbf{hard-fixed} at the ICBM value $k_1 = 0.8$ with a one-time reference-climate offset
($\alpha_A = k_1/\xi_{\mathrm{Ultuna}} = 0.851$); the slow rate stays free under a
\textbf{very-informative} prior centred on $k_2/\xi_{\mathrm{Ultuna}}$, on the argument that
$k_2$ is the one weak link (it is measured on Ultuna \emph{arable bare fallow}) and that the
posterior width would then serve as an arable$\rightarrow$forest transferability diagnostic.
Prior SD was tightened from $0.50$ to $0.15$ on the log scale.

\subsection{Humification and transfer fractions}
Free in every model, under the common Tier-2 logit prior with SD $0.4$, re-centred for the
simple models at the ICBM humification coefficient $h = 0.13$. Fixing them would have made the
simple models \emph{more} constrained than Yasso, which the intercomparison design forbids.

\subsection{The two auxiliary litter parameters}
$\sinput$ multiplies litter inputs; $\sinit$ scales the litter flux at the start of the
$1917\rightarrow1985$ pre-run.

\paragraph{Why they are not sampled directly.} A prior placed on $\sinput$ cannot express a
physical belief, because $\sinput$ is a bare multiplier with no referent --- ``$\sinput = 3$'' is
neither plausible nor implausible until you know what it multiplies. What \emph{is} physically
meaningful is the resulting litter \textbf{flux}, which has a boreal NPP envelope. So the
\texttt{flux\_pair} transform makes the two \emph{fluxes} the sampled quantities and recovers the
multipliers from them.

\paragraph{How the constraint works.} The sampler proposes two ordinary unconstrained real
numbers $x_1, x_2 \in (-\infty,\infty)$ --- it never sees a bound and never rejects a proposal.
Each is passed through the logistic, which maps the whole real line into $(0,1)$, and the result
is rescaled onto the physical window $(a,b)$:
\begin{equation}
F \;=\; a + (b-a)\,\frac{1}{1+e^{-x}} \;\in\; (a,b) \quad\text{for every } x\in\mathbb{R}.
\label{eq:squash}
\end{equation}
The bound is therefore a property of the \emph{coordinate system}, not a rule enforced during
sampling: an out-of-envelope flux is not rejected, it is unreachable. As $x\to\pm\infty$ the flux
approaches $a$ or $b$ asymptotically but never attains them. A Jacobian term
$\log|\mathrm{d}F/\mathrm{d}x|$ is added to the posterior so that this stretching does not
silently distort the prior.

\paragraph{The pair.} Applying eq.~\eqref{eq:squash} to both coordinates and recovering the
model parameters:
\begin{align}
F_{\mathrm{now}}  &= a + (b-a)\,\mathrm{logit}^{-1}(x_1), &
\sinput &= F_{\mathrm{now}}/\bar{J}, \label{eq:fp1}\\
F_{1917} &= a + (b-a)\,\mathrm{logit}^{-1}(x_2), &
\sinit  &= F_{1917}/F_{\mathrm{now}}, \label{eq:fp2}
\end{align}
with $[a,b] = [0.05,\,8.7]$\,tC\,ha$^{-1}$\,yr$^{-1}$ and $\bar{J} = 2.48$\,tC\,ha$^{-1}$\,yr$^{-1}$
the cross-plot mean litter, which converts an absolute flux into a dimensionless multiplier. It
is called a \emph{pair} because the two are coupled: $\sinit$ is not sampled at all but emerges
as the ratio of the two fluxes, and therefore depends on $x_1$ as well as $x_2$.

\paragraph{Worked example (Yasso20, current posterior).} $\sinput = 1.388$ implies
$F_{\mathrm{now}} = 1.388 \times 2.48 = 3.44$; $\sinit = 1.236$ implies
$F_{1917} = 1.236 \times 3.44 = 4.26$. Both sit inside $[0.05,\,8.7]$, so nothing in the
transform objects --- even though the second says 1917 forests produced 24\,\% more litter than
the contemporary flux. This is the concrete form of defect D3 below.

\begin{caveat}
\textbf{Two different ratios appear in this document and they are not the same number.}
$\sinit = F_{1917}/F_{\mathrm{now}} = 1.24$ compares the 1917 flux to the \emph{contemporary}
one. The quantity that decides whether the pre-run rises or falls is
$J_{1917}/J_{1985} = \sinit/0.818 = 1.51$, which compares it to the \emph{1985} flux. They differ
because the two ends of the pre-run are referenced to different litter aggregates --- which is
defect D1, and the reason this parameter is so easy to misread.
\end{caveat}

\paragraph{Where the priors actually live.} They are normal with SD $0.50$ on the
\emph{unconstrained} coordinates $x_1, x_2$, centred at the values that map to $\sinput = 1.30$
(C4b, raised above~1 to stand in for the missing understorey component) and $\sinit = 1.0$. A
consequence worth stating: because the prior sits on $x$ rather than on the flux, its shape in
physical units is not lognormal or anything else nameable --- it is a logit-normal on the window,
so its spread in tC\,ha$^{-1}$\,yr$^{-1}$ depends on where in the window it sits. This is why
neither of us could say what the $\sinit$ prior \emph{was} until we read the transform code.

\paragraph{It fixed what it was built for.} With a weak prior on the abstract multiplier, TP2 and
TP3 had driven $\sinput$ to $13$--$20\times$, implying effective litter fluxes of $34$ and
$50$\,tC\,ha$^{-1}$\,yr$^{-1}$ --- above boreal NPP. After \texttt{flux\_pair}, effective fluxes
are $0.9$--$3.5$ in all six models. That pathology is genuinely gone, which is also why the hard
bound may no longer be load-bearing (section~\ref{sec:package}, decision 3).

\subsection{Transient initialisation}
Each plot is started at steady state under a 1917 flux and the flux is then ramped to its 1985
value over 68 years:
\begin{equation}
J_{1917} = \Jfull \,\sinit\, \sinput, \qquad
J_{1985} = \Jto \, \sinput,
\label{eq:anchors}
\end{equation}
where $\Jfull$ is the plot's mean litter over the \emph{whole} 1986--2024 series and $\Jto$ its
mean over the \emph{first five} years. The ramp between them follows the national growing-stock
trajectory (C3, Korhonen et al.\ 2024), which replaced a linear ramp; the source file states
explicitly that ``litter scales with standing biomass'' and equally explicitly that ``the
calibrated endpoints $J_{1917}/J_{1985}$ are UNCHANGED; only the shape between moves.''

\subsection{Likelihood}
Multiplicative normal, with the SD proportional to the prediction and its coefficient fixed:
\begin{equation}
\mathrm{SOC}^{\mathrm{obs}} \sim \mathcal{N}\!\left(\widehat{\mathrm{SOC}},\;
\left(\widehat{\mathrm{SOC}}\cdot\sobs\right)^{2}\right), \qquad \sobs = 0.4416 .
\label{eq:lik}
\end{equation}
$\sobs$ is estimated externally from the observation campaigns and held fixed. There is no
separate model-error term.

% =====================================================================
\section{Defects identified 2026-08-07}\label{sec:defects}

\begin{defect}{D1 --- the two ends of the pre-run use different litter aggregates}
In eq.~\eqref{eq:anchors} the 1917 flux is referenced to $\Jfull$ (the 1986--2024 mean) and the
1985 flux to $\Jto$ (the 1986--1990 mean). Since litter rises through the record,
$\Jto/\Jfull$ has median $0.818$ over the 512 plots. Therefore
\[
\frac{J_{1917}}{J_{1985}} \;=\; \sinit \cdot \frac{\Jfull}{\Jto} \;=\; \frac{\sinit}{0.818},
\]
so $\sinit$ is \emph{not} the 1917/1985 flux ratio: it is that ratio scaled by an arbitrary
choice of reference aggregate. $\sinit = 1$ does not mean ``1917 like 1985''; it means ``1917
flux 22\,\% above 1985''. \texttt{CLAUDE.md} additionally still describes $\sinit$ as a
``proportional SD on the initial carbon state (log-transformed)'', which was true before
\texttt{flux\_pair} and is now wrong in both respects.
\end{defect}

\begin{defect}{D2 --- the initialisation prior is informative in the wrong direction}
The default $\sinit = 1.0$ therefore asserts a 1917 litter flux 22\,\% \emph{above} 1985.
The growing-stock record --- the same source used for the ramp shape --- implies a ratio below
one (section~\ref{sec:R}). The prior centre is roughly $1.5\times$ too high on the quantity
that sets the initial carbon state. This is not a weak prior that the data can overrule: the
posterior $\sinit$ moves $+0.00\sigma$ (Yasso15) to $+0.42\sigma$ (Yasso20) from that centre.
\end{defect}

\begin{defect}{D3 --- the constraint bounds levels, not the ratio}
\texttt{flux\_pair} bounds $F_{\mathrm{now}}$ and $F_{1917}$ each to $[a,b]$, but the pre-run
\emph{direction} depends on their ratio, which can take any value in $[a/b,\,b/a] = [0.006,\,174]$
without approaching a boundary. In the current run Yasso20 has $F_{1917} = 4.26$ against
$F_{\mathrm{now}} = 3.44$ --- more litter in 1917 than today --- with both fluxes comfortably
inside the envelope. The constraint is orthogonal to the failure.

Consequences in the current run: the pre-run declines for 78\,\% (Yasso15) and 92\,\% (Yasso20)
of plots, those two models over-predict the 1985 stock by $+19.0$ and $+23.5$\,tC\,ha$^{-1}$,
and both report a soil carbon \emph{source} ($-0.063$, $-0.155$\,tC\,ha$^{-1}$\,yr$^{-1}$)
where the three campaigns show a \emph{sink} ($+0.312$). The sign error survives dropping the
1985 campaign entirely.
\end{defect}

\begin{defect}{D4 --- the likelihood conflates observation error with model error}
$\sobs$ in eq.~\eqref{eq:lik} is fixed at the observation CV, and no model-error term exists.
Actual residuals exceed the assumed SD by a factor $1.19$--$1.21$ in \emph{all six} models.
Because the SD is proportional to the prediction, the only way the fit can widen its error
bars is to raise the predictions.
\end{defect}

% =====================================================================
\section{Proposed package --- NOT IMPLEMENTED}\label{sec:package}

\begin{caveat}
\textbf{Nothing in this section is in the code.} Section~\ref{sec:current} is what runs today;
this is what we propose to replace it with. the common anchor is a defect correction we would make regardless;
the ratio coordinate and the ratio prior are conditional on the unresolved level offset (section~\ref{sec:open}) --- see
section~\ref{sec:effect} for why they make two models visibly worse until that is understood.
\end{caveat}

The three proposals below are stated as a package because they are not independent: the ratio coordinate is
what makes the ratio prior expressible, and the common anchor is what makes the ratio coordinate mean what it says.

\begin{proposal}{COMMON ANCHOR --- reference both ends of the pre-run to the same aggregate}
Replace eq.~\eqref{eq:anchors} with
\begin{equation}
J_{1985} = \Jto\,\sinput, \qquad J_{1917} = \Jto\,\sinput\,R ,
\label{eq:newanchor}
\end{equation}
so that $R$ \emph{is} the 1917/1985 litter ratio. $\Jfull$ ceases to be used as a historical
anchor, which it should never have been: it is a modern-record mean standing in for a
historical quantity. $\sinit$ is retired and replaced by $R$.

\smallskip
\emph{Justification.} This is a defect correction, not a modelling choice. It should be made
regardless of what is decided about priors or bounds.
\end{proposal}

\begin{proposal}{RATIO COORDINATE --- make the second \texttt{flux\_pair} coordinate the ratio}
Keep the NPP envelope on the contemporary flux; sample the ratio rather than the 1917 flux:
\begin{align*}
F_{\mathrm{now}} &= a + (b-a)\,\mathrm{logit}^{-1}(x_1), & \sinput &= F_{\mathrm{now}}/\bar{J},\\
R &= R_{\mathrm{lo}} + (R_{\mathrm{hi}}-R_{\mathrm{lo}})\,\mathrm{logit}^{-1}(x_2), &
J_{1917} &= \Jto\,\sinput\,R,
\end{align*}
with $[R_{\mathrm{lo}}, R_{\mathrm{hi}}] \approx [0.25,\,1.25]$. The Jacobian gains a bounded
term for $R$ and loses the $-\log F_{\mathrm{now}}$ term, since the ratio no longer depends on
the contemporary flux. A guard returns $-\infty$ if $J_{1917}$ leaves the NPP envelope; with
$F_{\mathrm{now}} \approx 3.2$ this is far from binding, but it preserves the physical
guarantee explicitly rather than by assumption.

\smallskip
\emph{Justification.} The pre-run direction is the quantity that fails, and it is the quantity
for which external evidence exists. Placing the prior on it is the only way to state that prior
in physical units.
\end{proposal}

\begin{proposal}{RATIO PRIOR --- derive the prior on $R$ from the growing-stock record}
$R \sim \mathrm{Lognormal}(\log 0.90,\ 0.15\text{--}0.20)$, derived in section~\ref{sec:R}.

\smallskip
\emph{Justification.} It uses the same NFI source, and the same interpolation, that already
sets the ramp shape --- removing the inconsistency in which the pipeline invokes ``litter scales
with standing biomass'' for the shape and abandons it for the level.
\end{proposal}

\begin{caveat}
\textbf{This prior will do real work, and that must be stated.} The likelihood is nearly flat in
this direction: a 20\,\% error in the 1985 stock is half of $\sobs$, and the C5 sweep moved
$\sinit$ by a factor of four while changing trusted-campaign RMSE by $0.45\,\%$. A prior that
dominates a direction the data cannot see is not a neutral choice.

The defence is not that the prior is weak, but that \emph{the status quo also dominates that
direction} --- asserting $R \approx 1.22$ --- and does so in contradiction with evidence the
same pipeline already uses. Replacing an unjustified informative prior with a justified one
removes an inconsistency rather than adding an assumption. The sensitivity should nevertheless
be reported, as the C5 sensitivity now is.
\end{caveat}

\subsection*{plain priors --- can \texttt{flux\_pair} simply be removed?}

the common anchor to the ratio prior keep the bounded transform and change what the second coordinate means. A simpler
question is whether the transform is needed at all, since its cost is precisely that its priors
cannot be stated in physical units (section~\ref{sec:current}).

\paragraph{Why it might not be.} \texttt{flux\_pair} was introduced against a genuine pathology:
under the old Tier-3 prior (log SD $0.50$) TP2/TP3 drove $\sinput$ to $13$--$20\times$. But an
SD-$0.50$ lognormal places $\approx 2.4\,\%$ of its mass above the NPP ceiling, whereas an
SD-$0.30$ lognormal places $\approx 0\,\%$. The containment may therefore be achievable with a
tighter \emph{ordinary} prior. Indeed the prior currently implied by the transform is nearly
indistinguishable from $\mathrm{Lognormal}(\log 1.30,\,0.30)$: quartiles $1.04/1.30/1.58$ against
$1.06/1.30/1.59$, differing only in a slightly thinner right tail ($2.14$ vs $2.34$ at 97.5\,\%).

\paragraph{Test.} \texttt{doublechecks/preflight\_naive\_priors.R} patches the real calibration
script to swap \texttt{flux\_pair} for two \texttt{log} parameters at
$\mathrm{SD} = 0.30$ / $0.20$, leaving everything else --- anchoring, centres, data --- untouched,
and pushes $K=200$ prior draws through the genuine \texttt{ll\_fn} at full $N$. Both
configurations use the same seed, so the comparison is paired.

\begin{table}[H]\centering\small
\caption{Prior-pushforward, $K=200$ draws at full $N=410$. ``rejects'' are legitimate
stick-breaking simplex rejections; ``blow-ups'' are non-finite $\ell$ among constraint-passing
draws --- the MCMC \texttt{inf\_rate} proxy.}
\begin{tabular}{lccc}
\toprule
model & naive (no \texttt{flux\_pair}) & production (\texttt{flux\_pair}) & rejects (both)\\
\midrule
SP1     & 0 / 200        & 0 / 200        & 0\\
TP2     & 0 / 200        & 0 / 200        & 0\\
TP3     & 0 / 200        & 0 / 200        & 0\\
Yasso07 & 5 / 101 (5.0\,\%) & 5 / 101 (5.0\,\%) & 99\\
Yasso15 & 0 / 145        & 0 / 145        & 55\\
Yasso20 & 0 / 145        & 0 / 145        & 55\\
\bottomrule
\end{tabular}
\end{table}

\noindent \textbf{Identical in every model.} Yasso07's five blow-ups occur in both configurations
among the same 101 constraint-passing draws, so they are the known Yasso07 $\delta_2/\sigma$
baseline rather than a consequence of removing the bound. Implied effective flux under the naive
prior is median $2.97$--$3.35$ with 97.5\,\% at $4.94$--$6.34$\,tC\,ha$^{-1}$\,yr$^{-1}$, and
\textbf{$0\,\%$ of prior mass exceeds the NPP ceiling in all six} --- the bound would have had
nothing to remove.

\begin{proposal}{PLAIN PRIORS --- the simplified parameterization}
\begin{align*}
\sinput &\sim \mathrm{Lognormal}(\log m,\ 0.30) & &\text{$m$: see decision 5}\\
R &\sim \mathrm{Lognormal}(\log 0.90,\ 0.20) & &\text{growing-stock elasticity, \S\ref{sec:R}}\\
J_{1985} &= \Jto\,\sinput, \qquad J_{1917} = \Jto\,\sinput\,R\\
\text{guard} &: \text{reject if the effective flux exceeds the NPP ceiling}
\end{align*}
Four statements, each expressible in one sentence of Methods, each carrying a prior on a
quantity with a referent. It subsumes the common anchor to the ratio prior and \emph{dissolves} D1 and D3 rather than
repairing them: there is no level-versus-ratio mismatch left to have.
\end{proposal}

\begin{caveat}
\textbf{This tests the prior, not the posterior.} The pre-flight confirms the forward model stays
finite across prior mass. It cannot rule out the sampler wandering into a large-$\sinput$ region
that the bound would have blocked --- a lognormal has an infinite tail where the logistic has a
wall, and the $13$--$20\times$ excursions were a \emph{posterior} phenomenon. Only a real
calibration settles that, which is why the guard is retained as a truncation that should never
fire. Removing \texttt{flux\_pair} would also require reframing the F9 result, currently stated
as ``the flux bound closed the $\sinput$ escape hatch''; the honest replacement is that the
escape hatch was a symptom of an inflated SOC target and a starved cascade, and that once those
were corrected an informative physical prior sufficed.
\end{caveat}

% =====================================================================
\section{Deriving \texorpdfstring{$R$}{R}}\label{sec:R}

\paragraph{Why not the volume ratio.} The NFI record gives growing stock $1400$\,Mm$^3$ at 1917
(NFI1, 1922, held flat back over a near-stationary period) and $1775$\,Mm$^3$ at 1985, a ratio
of $0.789$. Taking $R = 0.789$ assumes litter is \emph{proportional} to standing volume. Three
arguments say it is not: foliage and fine-root litter scale sub-linearly with stem volume
because leaf area saturates as stands close; woody litter tracks mortality rather than standing
stock, and the mortality regime of cut-over unmanaged forest differs; and understorey litter ---
absent from the litter product entirely --- is proportionally \emph{larger} in sparse stands.
All three bias $R$ upward relative to the volume ratio.

\paragraph{Empirical elasticity.} Rather than assume an exponent, it can be fitted. Both series
exist over 1986--2023: national mean litter from the input bundle, and growing stock
interpolated from the same NFI file used for the ramp shape. Fitting $J \propto \mathrm{GS}^{\epsilon}$
and extrapolating $R = (\mathrm{GS}_{1917}/\mathrm{GS}_{1985})^{\epsilon}$:

\begin{table}[H]\centering\small
\caption{Litter--growing-stock elasticity over 1986--2023 (growing stock $+42\,\%$, mean litter
$+35\,\%$ over the window), and the implied 1917/1985 litter ratio.}
\begin{tabular}{lrrr}
\toprule
fit & $\epsilon$ & $R = 0.789^{\epsilon}$ & $R^2$\\
\midrule
annual, all years          & 0.467 & 0.895 & 0.370\\
5-year block means         & 0.448 & 0.899 & 0.401\\
endpoints only (5-yr means)& 0.661 & 0.855 & ---\\
\midrule
\emph{naive proportionality} & \emph{1.000} & \emph{0.789} & ---\\
\bottomrule
\end{tabular}
\end{table}

\noindent The elasticity is close to $0.5$ --- litter scales roughly with the square root of
growing stock --- and is stable against interannual climate noise, which the low $R^2$ of the
annual fit makes the obvious concern. The regression SE on $\epsilon$ ($0.101$) alone would
give $R \in [0.854, 0.938]$.

\paragraph{Adopted value.} We propose centring at $R = 0.90$ with log SD $0.15$--$0.20$, wider
than the regression SE alone would justify, to accommodate the caveats below.

\begin{caveat}
\textbf{Five caveats on this derivation.}
(i) \emph{Possible circularity} --- if the litter product is itself derived from biomass, this
elasticity is internal to the product rather than independent evidence. For the purpose of
extrapolating \emph{that product} backwards consistently it is arguably the correct choice, but
it requires confirmation from B.~Tupek on how the product is constructed. This is the caveat
most likely to matter to a reviewer.
(ii) \emph{Extrapolation} --- the relation is fitted over GS $1798$--$2552$ and applied at
$1400$, $22\,\%$ below the fitted range.
(iii) \emph{Understorey omitted} --- the product is tree litter only, so the total-litter ratio
is biased \emph{upward} from this estimate, plausibly $0.92$--$0.95$. Our value is a lower bound.
(iv) \emph{A national ratio applied per plot} --- individual plot histories differ from the
national mean.
(v) \emph{Non-stationarity} --- management, species composition and age structure changed over
1917--1985, so the modern volume-to-litter relation need not have held historically.
\end{caveat}

% =====================================================================
\section{What the package fixes, and what it does not}\label{sec:effect}

\paragraph{Fixes.} $\sinit$ is retired in favour of a parameter with a referent and a
stateable prior; the pre-run direction becomes evidence-constrained; the shape/level
inconsistency in the C3 implementation disappears; and --- if the mechanism holds --- the
Yasso15/20 pre-run inversion and their inverted sink go with it.

\paragraph{It does not fix the level.} A uniform over-prediction of $17$--$20\,\%$ is shared by
all six models (section~\ref{sec:open}), and the common anchor to the ratio prior do not address it.

\paragraph{And it makes two models visibly worse.} At $R = 0.90$ the equivalent $\sinit$ is
$\approx 0.74$. Against the current posteriors:

\begin{table}[H]\centering\small
\caption{Current $\sinit$ against the proposed centre (displacement at log SD $0.20$).}
\begin{tabular}{lrr}
\toprule
model & $\sinit$ & vs.\ proposed centre\\
\midrule
SP1     & 0.627 & $-0.8\sigma$\\
TP2     & 0.564 & $-1.4\sigma$ \emph{(too low)}\\
TP3     & 0.566 & $-1.4\sigma$ \emph{(too low)}\\
Yasso07 & 0.768 & $+0.2\sigma$\\
Yasso15 & 1.002 & $+1.5\sigma$\\
Yasso20 & 1.236 & $+2.6\sigma$ \emph{(too high)}\\
\bottomrule
\end{tabular}
\end{table}

\noindent The constraint pulls Yasso15/20 down, which is the intended repair, but pulls TP2/TP3
\emph{up}, giving them more 1917 input and worsening their level over-prediction. TP2 and TP3
have evidently been suppressing their 1917 input below what the forest history supports,
partially offsetting the level bias. Pinning $R$ removes that compensation and forces the
misfit into the open --- the same logic by which the \texttt{flux\_pair} bound became a result
rather than a cosmetic fix. It does, however, mean the two problems cannot be adjudicated
independently: if $R$ is pinned before the level offset is understood, the simple models will
get visibly worse and it will not be clear whether that is the constraint working as intended
or a second error compounding.

% =====================================================================
\section{Open: the uniform level offset}\label{sec:open}

All six models over-predict by $17$--$20\,\%$ (median log-residual $-0.176$ to $-0.206$;
campaign means $+22$--$38\,\%$ at 1985 falling to $+7$--$14\,\%$ at 2024). Three explanations
were tested on 2026-08-07 and none is fully satisfactory.

\begin{itemize}
\item \emph{Not the input data.} $\sinput$ is a free per-model multiplier on exactly that
quantity and demonstrably moves (SP1 sits $2.6\sigma$ from its prior centre; no flux bound
binds). Any uniform input bias is absorbed by construction.
\item \emph{Not the $\sinput$ prior centre.} SP1 is $2.6\sigma$ \emph{below} the C4b centre and
still over-predicts by $24\,\%$; and the excess is $23$--$28\,\%$ in all six models while
$\sinput$ itself spans $0.35$--$1.39$.
\item \emph{Not insufficient kinetic freedom.} The rates moved $1.6$--$2.2\sigma$ against a
prior of SD $0.15$ --- but \emph{toward longer} residence times ($+28$ to $+38\,\%$), i.e.\
\emph{raising} SOC. Both available routes to lowering the level were left unused or used in the
opposite direction.
\item \emph{Partly the likelihood form.} Under eq.~\eqref{eq:lik} the optimal uniform rescaling
of the predictions is $k = 1.019$: the likelihood positively prefers the current
over-prediction. Freeing $\sobs$ within the same form moves this only to $k = 0.960$; changing
the form to log-normal moves it to $k = 0.848$, against $0.811$ for matching the observed
median. So the \emph{form} accounts for most of the offset and the \emph{width} for little ---
but this is inference from fitted predictions, not a refit.
\end{itemize}

\begin{caveat}
The decisive test has not been run: refit one model (TP2) under a log-normal likelihood with
nothing else changed. If the offset collapses and $\sinput$ lands near the $1.04$ the
observations require, the fitting criterion is responsible. If it does not, the cause is
neither the inputs, the priors, the kinetics nor the likelihood, and remains unidentified.
\end{caveat}

\paragraph{A note on the rate displacement.} C1's design states that the slow-rate posterior
width would serve as an arable$\rightarrow$forest transferability diagnostic. That diagnostic
has fired: all three simple models sit $1.6$--$2.2\sigma$ from the ICBM anchor, consistently
slower, implying Finnish forest slow-pool turnover some $30\,\%$ slower than the Ultuna arable
value. Whether this is a genuine transferability result or the rates absorbing misfit that
belongs elsewhere is testable --- if the latter, correcting the level should pull them back
toward the anchor.

% =====================================================================
\section{The level and the trend are different problems}\label{sec:leveltrend}
% =====================================================================

Section~\ref{sec:open} recorded a uniform $\sim$20\,\% over-prediction of unknown cause. It has
since been identified, tested and largely removed --- and in doing so it separated cleanly from a
second problem that it had been masking.

\subsection{The level: an error-model artefact, now fixed}

The multiplicative normal of eq.~\eqref{eq:lik} sets the variance from the parameter being
estimated, so a prediction widens its own error bar. The penalty for a badly-missed plot
therefore \emph{plateaus} rather than growing, and the fit buys tolerance for its worst plots by
inflating everything. The consequence is a location estimate contaminated by dispersion. On data
generated \emph{unbiased} by construction, the optimal scaling under this form drifts from $0.86$
to $3.25$ as the residual CV rises from $0.05$ to $1.0$, while a log-normal returns $1.000$ at
every spread.

A refit of TP2 with the likelihood replaced by a log-normal (3 chains $\times$ 6000, C5 $=1.0$;
control is the matched \texttt{A1\_C5\_off} ablation) confirms it:

\begin{table}[H]\centering\small
\caption{TP2 under the two error models, at posterior-median parameters.}
\begin{tabular}{lrr}
\toprule
 & multiplicative normal & log-normal\\
\midrule
bias 1985 (tC\,ha$^{-1}$) & $+15.2$ & $+2.8$\\
bias 2006 & $+11.6$ & $-2.2$\\
bias 2024 & $+11.5$ & $-3.1$\\
\textbf{bias overall} & $\mathbf{+12.7}$ & $\mathbf{-0.9}$\\
median log-residual & $-0.205$ & $-0.023$\\
\midrule
paired trend 1985--2024 & $+0.168$ & $+0.116$\\
\emph{observed} & \multicolumn{2}{c}{$+0.309$}\\
\bottomrule
\end{tabular}
\end{table}

\noindent The bias collapses to essentially zero. The implied rescaling, $0.834$, matches the
$0.848$ predicted from the closed-form analysis to within 2\,\%. Convergence is acceptable
(R-hat $1.009$--$1.033$, ESS $206$--$382$). \textbf{The level bias was ours, not the models'} ---
which means every level-based number from the \texttt{20260805} run carries it.

\begin{caveat}
The route was not the one predicted. $\sinput$ barely moved ($1.290 \rightarrow 1.306$); the model
took the level down through the climate parameters instead ($\beta_1$ $+35\,\%$). The
\emph{magnitude} of a refit is predictable from a fixed-prediction rescaling analysis; the
\emph{parameter that delivers it} is not. Relatedly, $\sinit$ rose ($0.563 \rightarrow 0.622$,
$R$ $0.69 \rightarrow 0.76$), so part of the apparent ``the data already infers historical
depletion'' signal was the biased likelihood.
\end{caveat}

\subsection{The trend: a response-time problem, with two distinct causes}

Correcting the level did not fix the trend --- it slightly worsened it, and changed its
signature from a uniform offset to the model \emph{crossing} the observations. That is expected:
a trend is governed by how far below its own equilibrium a model starts, and that ratio is
invariant under uniform rescaling.

Fitting an exponential approach $C(t) = E - (E-C_0)e^{-t/\tau}$ to each model's own trajectory
(1985 / 2024 / 2084) gives its effective response time:

\begin{table}[H]\centering\small
\caption{Effective response time and starting position, fitted per model. The observations
require $\tau \approx 44$--$87$\,yr and a start at $\approx 0.73$ of equilibrium.}
\begin{tabular}{lrrrr}
\toprule
model & $\tau$ (yr) & start/eq & gap closed in 39\,yr & 2024 if re-initialised at 61.5\\
\midrule
SP1     & 79  & 0.79 & 39\,\% & 75.4 \quad($+0.357$)\\
TP2     & \textbf{195} & 0.68 & \textbf{18\,\%} & 70.8 \quad($+0.237$)\\
TP3     & \textbf{194} & 0.68 & \textbf{18\,\%} & 70.9 \quad($+0.240$)\\
Yasso07 & 38  & 0.93 & 64\,\% & 76.3 \quad($+0.381$)\\
Yasso15 & \textbf{17}  & \textbf{1.03} & 90\,\% & 77.6 \quad($+0.412$)\\
Yasso20 & 23  & \textbf{1.08} & 82\,\% & 76.3 \quad($+0.379$)\\
\midrule
\emph{observed} & \emph{44--87} & \emph{0.73} & \emph{36--57\,\%} & \emph{75.4 \quad($+0.309$)}\\
\bottomrule
\end{tabular}
\end{table}

\noindent The failure has \textbf{two different causes}:

\paragraph{Yasso (and SP1): an initialisation problem.} Their response times are already at or
beyond the required range. All three Yasso models start at or \emph{above} their own equilibrium
(start/eq $0.93$, $1.03$, $1.08$) and therefore decline or stagnate. Re-initialised at the
observed 1985 stock, keeping each model's own $E$ and $\tau$, the three reach $+0.381$, $+0.412$
and $+0.379$ against an observed $+0.309$, and SP1 lands on $75.4$ exactly. For this family the
initialisation constraint (the ratio prior) is the operative fix and the rate anchor is irrelevant --- Yasso's
rates are fixed by design.

Note however that the Yasso response times ($17$--$38$\,yr) are \emph{faster} than the
$44$--$87$\,yr the observations imply, so correcting the initialisation alone would leave them
over-accumulating by $23$--$33\,\%$: they would swing from reporting a source to over-reporting a
sink. That is a large improvement on a sign error, and the counterfactual holds $E$ fixed whereas
a real recalibration would adjust it --- but ``necessary and sufficient'' overstates it. Necessary,
and sufficient for the sign and the order of magnitude.

\paragraph{TP2/TP3: a turnover problem.} At $\tau = 195$\,yr they close only 18\,\% of the gap
and reach $70.8$ even when perfectly initialised. No initialisation fixes them. They require
$\alpha_H$ some $2.3$--$4.7\times$ faster than the ICBM anchor, and --- because faster turnover
also lowers the equilibrium --- the increase must be accompanied by a compensating rise in
$\sinput$ so that the level is held while the response speeds up.

\begin{proposal}{FREE SLOW RATE --- relax the slow-rate anchor for the simple models}
Widen the C1 slow-rate prior (currently log SD $0.15$, centred on Ultuna arable $k_2$) so the
data can select the turnover time. Report the displacement as the transferability result.

\smallskip
\emph{Why this is not a retreat from C1.} C1's stated design makes the slow-rate posterior width
``the arable$\rightarrow$forest transferability diagnostic''. It has fired: all three simple
models sit $1.6$--$2.2\sigma$ from the anchor, consistently slower, and the trend analysis says
the required correction is $2.3$--$4.7\times$. The \emph{fast} rate stays hard-fixed --- it is
litterbag-grounded and transferable; only the component whose provenance was flagged as weak
(bare fallow, no fresh input, most recalcitrant fraction remaining) is freed.

\smallskip
\emph{It corrects an asymmetry rather than creating one.} The homogenisation principle is equal
external \emph{information}. Yasso's rates are fixed because they were calibrated on forest data;
the simple models' were anchored to arable data. Those are not of equal quality for a forest
application, and freeing the latter restores the intended parity.

\smallskip
\emph{Independent corroboration.} Yasso's forest-calibrated rates give effective
$\tau = 23$--$38$\,yr, against $195$\,yr for the ICBM-anchored cascades. Two independent
sources --- Yasso's own calibration and the observed Finnish SOC trend --- agree that forest
turnover is several times faster than the Ultuna arable value.
\end{proposal}

\begin{caveat}
These response times are fitted from three trajectory points using the 2084 projection as an
equilibrium proxy, on cross-plot means. Directionally solid; the magnitudes need an analytical
treatment before publication. And the observed target is itself uncertain: on the two campaigns
validated against official LUKE figures (2006$\rightarrow$2024) the observed rate is $+0.209$,
not $+0.309$, so part of the apparent deficit may be the 1985 stock reading low --- the C5
concern resurfacing on the trend rather than the level.
\end{caveat}

\subsection{Can the simple structures fit at all?}\label{sec:cansimplefit}

The response-time analysis appears to condemn TP2/TP3: at $\tau = 195$\,yr they close only
18\,\% of the gap and cannot reach the observed accumulation at any initial state. But that
conclusion depends on which parameters are held fixed. TP2 has \emph{three} relevant free knobs,
not one:
\[
E = J_{\mathrm{eff}}\left[\frac{1}{\alpha_A \xi} + \frac{p_H}{\alpha_H \xi}\right],
\qquad \tau = \frac{1}{\alpha_H \xi},
\]
so $\alpha_H$ sets the response time, $p_H$ raises the equilibrium \emph{without} slowing the
response, and $\sinput$ scales both. Faster turnover alone lowers $E$ proportionally and is
self-defeating; faster turnover \emph{with} compensating $p_H$ or $\sinput$ is not.

\begin{table}[H]\centering\small
\caption{TP2 feasibility at $\tau = 45$\,yr (the fast end, best matching the observed
deceleration), for the equilibrium $E = 85.5$ implied by $61.5 \rightarrow 75.4$ over 39\,yr.
NPP ceiling $8.7$\,tC\,ha$^{-1}$\,yr$^{-1}$; $p_H$ displacement is against the Tier-2 prior
(centre $0.13$, logit SD $0.4$).}
\begin{tabular}{rrrrl}
\toprule
$p_H$ & $\sinput$ & flux & $p_H$ displacement & \\
\midrule
0.13 (ICBM) & 4.81 & 11.94 & $+0.0\sigma$ & \textbf{exceeds NPP}\\
0.20        & 3.34 & 8.29  & $+1.3\sigma$ & ok, tight\\
\textbf{0.30} & \textbf{2.33} & \textbf{5.77} & $+2.6\sigma$ & \textbf{comfortable}\\
0.45        & 1.60 & 3.97  & $+4.3\sigma$ & very comfortable\\
\bottomrule
\end{tabular}
\end{table}

\noindent \textbf{So the simple structures can fit.} TP2 reaches the observed level \emph{and}
trend with $\alpha_H \approx 0.025$, $p_H \approx 0.30$ and $\sinput \approx 2.3$ --- an effective
flux of $5.8$, comfortably inside the envelope. The apparent impossibility was an artefact of
holding $p_H$ at its current value while varying $\alpha_H$.

\paragraph{What it costs relative to the arable analogue.} Humus turnover $\approx 3.9\times$
faster than Ultuna $k_2$; humification efficiency $0.13 \rightarrow 0.30$, a $2.6\sigma$ prior
displacement but squarely inside the range usually quoted for humification efficiencies
($0.1$--$0.3$); and litter input $2.3\times$ the measured product. Each is individually
defensible; together they are a substantial departure from ICBM, and the departure is the
interesting part --- particularly for $p_H$, where a forest organic layer humifying more
efficiently than arable bare fallow is close to expected.

\paragraph{Consequence for the intercomparison.} This inverts the reading of C1. The anchor was
not making the comparison fair; it was preventing the simple models from reaching the region of
parameter space that fits, and the $10\times$ response-time gap against Yasso was the symptom.
With forest-appropriate kinetics the two-pool cascade \emph{can} represent Finnish SOC dynamics,
which is a positive answer to the question the simple models were included to ask.

\begin{caveat}
This is a steady-state feasibility calculation on cross-plot means. It establishes that a
solution \emph{exists} in parameter space --- not that the sampler finds it, nor that it holds
plot-by-plot. The three knobs trade off along a ridge, so the calibration may land anywhere
along it; watch $\alpha_H$, $p_H$ and $\sinput$ \emph{together}, not $\alpha_H$ alone.
\end{caveat}

\subsection{Correction: the trend is driven by rising litter, not by historical depletion}
\label{sec:rising}

Sections~\ref{sec:leveltrend}--\ref{sec:cansimplefit} analysed the trend as relaxation toward a
\emph{fixed} equilibrium, and concluded that the models must start far below it --- with a
required $R \approx 0.60$, well under the growing-stock value, motivating a historical-depletion
hypothesis (litter raking, forest grazing, slash-and-burn). \textbf{That analysis was wrong and
the conclusion is withdrawn.}

The equilibrium is not fixed: national mean litter rises $\times 1.35$ over 1986--2024. A soil
starting \emph{at} equilibrium and simply tracking that rise produces:

\begin{table}[H]\centering\small
\caption{SOC in 2024 for a soil at equilibrium with 1985 litter, tracking the observed litter
rise with response time $\tau$. Observed: $75.4$, trend $+0.312$.}
\begin{tabular}{rrr}
\toprule
$\tau$ (yr) & SOC 2024 & trend\\
\midrule
23 (Yasso20)  & 76.6 & $+0.387$\\
30            & 74.8 & $+0.341$\\
45            & 71.9 & $+0.267$\\
195 (TP2)     & 64.6 & $+0.081$\\
\bottomrule
\end{tabular}
\end{table}

\noindent $\tau \approx 33$--$38$\,yr reproduces the observed accumulation exactly, with
\textbf{no below-equilibrium start required}. The diagnostic error was the same one that gave
SP1 a fitted $\tau$ of 79\,yr against its true value of 33: a three-point exponential fit
assumes a stationary target, and the target is not stationary.

\paragraph{Consequences.}
\begin{itemize}
\item The \textbf{historical-depletion hypothesis is not needed}. It may still be true, but the
SOC record does not demand it, and no claim should rest on it.
\item The \textbf{conflict over $R$ dissolves}. There is no competing trend requirement pulling
$R$ below the growing-stock value, so $R \approx 0.90$ (section~\ref{sec:R}) stands uncontested
and can be applied to all models without harming the cascades.
\item The \textbf{two failure modes are now cleanly separated}, one per family and sharing no
parameter: Yasso is initialised \emph{above} equilibrium ($R = 1.22$, $1.51$) and relaxes
downward --- fixed by $R$; TP2/TP3 are too slow to track the rise ($\tau = 195$) --- fixed by
$\alpha_H \approx 0.032$, roughly $4.9\times$ ICBM, and unrelated to $R$.
\end{itemize}

\paragraph{The mechanism, stated plainly.} Finnish forest soils are accumulating carbon because
litter input is rising with the growing stock. A model reproduces that only if its turnover is
fast enough to track the rise, and only if it is not initialised above its own equilibrium. Both
observed failure modes follow from that one sentence.

\subsection{$\sinit$ is conditionally identified, and the condition is the 1985 campaign}
\label{sec:cond1985}

It would be wrong to describe $\sinit$ as unidentified. Pushing the prior through the
\texttt{flux\_pair} transform and comparing with the posteriors:

\begin{table}[H]\centering\small
\caption{Implied prior on $\sinit$ versus the fitted posteriors (\texttt{20260805} run).}
\begin{tabular}{lrrr}
\toprule
 & median & 95\,\% width & narrowing\\
\midrule
\textbf{prior (implied)} & \textbf{1.000} & \textbf{1.950} & ---\\
SP1     & 0.627 & 0.322 & $6.1\times$\\
TP2     & 0.564 & 0.325 & $6.0\times$\\
TP3     & 0.566 & 0.332 & $5.9\times$\\
Yasso07 & 0.768 & 0.273 & $7.1\times$\\
Yasso15 & 1.002 & 0.397 & $4.9\times$\\
Yasso20 & 1.236 & 0.895 & $2.2\times$\\
\bottomrule
\end{tabular}
\end{table}

\noindent The data narrows $\sinit$ by $2$--$7\times$ and moves the median by up to $44\,\%$
(prior $1.00 \rightarrow$ TP2 $0.564$). Four of the six models pull it a long way down. This is a
parameter the calibration genuinely informs.

\paragraph{But the information sits almost entirely in one campaign.} 1985 is the observation
closest to the initial state, so it carries most of the signal about it. The C5 sweep is the
direct demonstration: changing how much the 1985 campaign is trusted moved $\sinit$ by a factor
of \emph{four} while trusted-campaign RMSE moved $0.45\,\%$. The parameter is well determined
\emph{given} a choice about 1985, and that choice is not itself determined by the data.

\paragraph{This is a scope statement, not a defect.} The correct form of the inference is
conditional: \emph{given that the VMI8 campaign is taken at face value, the reconstructed initial
state is $X$}; down-weighting VMI8 shifts it by up to $4\times$ at negligible cost in predictive
skill. Stating it that way is honest and it is exactly the sensitivity the C5 record was kept to
support. It should not be presented as an unconditional estimate.

It also \emph{localises} the initialisation problem usefully. The difficulty is not that soil
carbon data cannot in principle constrain an initial state --- it demonstrably can, by a factor
of $2$--$7$. It is that the constraint rests on the one campaign whose layer protocol differs
(VMI8 samples $0$--$5$ and $5$--$20$\,cm; the later campaigns $0$--$10$ and $10$--$20$\,cm) and
which could not be validated against official LUKE figures. That is a specific, addressable
data question rather than a generic identifiability lament.

\paragraph{And it explains Yasso20.} It is the \emph{least} constrained of the six --- $2.2\times$
narrowing against $6$--$7\times$ elsewhere, and a posterior nearly three times wider. So Yasso20
does not merely land in the wrong place; it is the model the data pins down least, which is why
it drifts furthest and why the prior matters most there. That is a sharper diagnosis of the
sink-sign failure than ``$\sinit$ is too high''.

\subsection{The initialisation cannot deliver the trend --- tested and rejected}
\label{sec:sweep}

Section~\ref{sec:leveltrend} attributed the Yasso trend failure to the initialisation. A direct
sweep of the 1917 anchor, at posterior-median parameters, shows that initialisation controls the
\emph{sign} but cannot deliver the \emph{magnitude}.

\begin{table}[H]\centering\small
\caption{Sweeping $\sinit$ at posterior-median parameters (log-normal fits). Observed:
1985 $=61.6$, 2024 $=73.8$, trend $+0.312$.}
\begin{tabular}{lrrr|rrr}
\toprule
 & \multicolumn{3}{c|}{TP2} & \multicolumn{3}{c}{Yasso20}\\
$\sinit$ & 1985 & 2024 & trend & 1985 & 2024 & trend\\
\midrule
0.10 & 15.6 & 29.5 & $+0.353$ & 16.9 & 26.4 & $+0.231$\\
0.25 & 26.8 & 38.5 & $+0.299$ & 23.3 & 31.3 & $+0.194$\\
0.55 & 49.1 & 56.4 & $+0.191$ & 36.0 & 41.0 & $+0.119$\\
0.70 & \textbf{60.3} & 65.4 & $+0.137$ & 42.4 & 45.8 & $+0.082$\\
1.10 & 90.2 & 89.3 & $-0.008$ & \textbf{59.4} & 58.7 & $-0.018$\\
\bottomrule
\end{tabular}
\end{table}

\noindent \textbf{The level and the trend make opposite demands on the anchor.} Matching the 1985
stock needs $\sinit \approx 0.70$ (TP2) or $\approx 1.15$ (Yasso20), which gives trends of
$+0.137$ and $-0.02$. Matching the trend needs anchors that put the 1985 stock at $27$ and
$17$\,tC\,ha$^{-1}$ respectively --- less than half the observed. \textbf{No anchor value satisfies
both, in either family.}

This rejects the historical-depletion hypothesis (litter raking, forest grazing, slash-and-burn)
as an explanation for the trend: it is implementable, and it fails. The hypothesis may remain true
as history; it is not the mechanism behind the trend deficit. What survives is that
\emph{initialisation controls the sink sign} --- Yasso20 spans $-0.09$ to $+0.23$ across the sweep
--- but not its magnitude.

\subsection{The trusted-campaign comparison cannot adjudicate 1985}
\label{sec:power}

A natural test of ``VMI8 reads low'' is to compare models over 2006--2024, which excludes it. That
test has almost no power:

\begin{table}[H]\centering\small
\caption{Observed paired rates and what a deficit over 2006--2024 could be detected.}
\begin{tabular}{lrrrr}
\toprule
period & $n$ & rate & SE & rate/SE\\
\midrule
1985--2006 & 348 & $+0.458$ & 0.048 & 9.5\\
2006--2024 & 409 & $+0.209$ & 0.065 & 3.2\\
1985--2024 & 316 & $+0.312$ & 0.035 & 9.0\\
\midrule
\multicolumn{5}{l}{\footnotesize\emph{Basis (added 2026-08-12): paired, unweighted, whole
profile --- correct for model comparison. Balanced (same plots as 1985--2024) gives $+0.143$;
region-weighted gives $+0.161$; LUKE's official weighted 0--40\,cm figure is $+0.111$. See
\texttt{doublechecks/observed\_soc\_basis.R}. Do not quote a weighted national stock in the
same sentence as a model--observation gap.}}\\
\bottomrule
\end{tabular}
\quad
\begin{tabular}{lr l}
\toprule
model captures & gap & \\
\midrule
30\,\% & $2.2$ SE & detectable\\
50\,\% & $1.6$ SE & \emph{indistinguishable}\\
70\,\% & $1.0$ SE & \emph{indistinguishable}\\
90\,\% & $0.3$ SE & \emph{indistinguishable}\\
\bottomrule
\end{tabular}
\end{table}

\noindent Over 2006--2024 the total change is $+3.8$\,tC\,ha$^{-1}$ on a stock of $\sim$70 --- a
5\,\% change at $3.2\sigma$ --- so a model capturing only half the trend is statistically
indistinguishable from one capturing all of it. Essentially all the discriminating information
lives in the 1985--2006 leg ($+9.6$\,tC\,ha$^{-1}$ at $9.5\sigma$), which is the leg whose
reliability is in question.

\textbf{So ``is 1985 low, or are the cascades too slow?'' cannot be resolved with these three
campaigns.} Both hypotheses predict the same observable, and the only period with power is the
contested one. This is a limit of the data, not an open analysis task.

(Yasso20 remains an exception: it predicts $-0.147$ against $+0.209$, a gap of $5.5$\,SE, so its
failure on the trusted campaigns stands regardless of 1985.)

\subsection{The target itself is not yet defined}\label{sec:target}

Every model--data gap above is quoted against an observed rate that varies substantially with two
choices we have not made explicitly.

\begin{table}[H]\centering\small
\caption{Observed paired rate, plot set defined by the data alone.}
\begin{tabular}{lrrrr}
\toprule
period & $n$ & mean & 10\,\% trimmed & median\\
\midrule
1985--2006 & 348 & $+0.458$ & $+0.409$ & $+0.410$\\
2006--2024 & 409 & $+0.209$ & $+0.105$ & $+0.071$\\
1985--2024 & 316 & $+0.312$ & $+0.260$ & $+0.229$\\
\bottomrule
\end{tabular}
\end{table}

\noindent \textbf{Statistic:} the 1985--2024 rate spans $+0.229$ (median) to $+0.312$ (mean), a
36\,\% range; over 2006--2024 the spread is a factor of three. The distribution of per-plot
changes is right-skewed.

\textbf{Plot set:} an earlier draft quoted $n = 328$ for 2006--2024, taken from plots where the
forward model happened to run. That is not a legitimate target --- it lets the model influence
what it is judged against. The data-defined set is $n = 409$.

The target therefore spans roughly $+0.07$ to $+0.31$ across defensible choices, which is the same
magnitude as the model deficit being explained. \textbf{One plot set and one statistic must be
fixed, and applied identically to observed and modelled rates}, before any gap is attributed to a
mechanism. For an inventory the mean is the defensible choice --- total stock change is what is
reported, and the median discards the high-accumulation plots that carry the sink --- but the skew
should be reported alongside, since a mean--median gap this large is itself a result about the
heterogeneity of accumulation.

\subsection{Acceptance criteria agreed for the next run}

\begin{table}[H]\centering\small
\begin{tabular}{ll}
\toprule
Yasso, effective flux (F9) & must stay well inside the NPP envelope --- non-negotiable\\
SP1/TP2/TP3, effective flux & high flux tolerable; report as a finding\\
Trend, all models & ``reasonably close'': within roughly a third of observed\\
Level, all models & aligned at 1985 (achieved by the likelihood fix)\\
\bottomrule
\end{tabular}
\end{table}

\noindent The two families need qualitatively different fixes, and only one strains plausibility:
\textbf{Yasso's correction is flux-neutral} (the R prior moves $\sinit$, not $\sinput$, so its
contemporary flux of $2.4$--$3.5$ is untouched), whereas \textbf{the cascades' correction costs
input} ($\sinput \approx 2.3$). For a study asking whether simpler structures suffice, that
asymmetry is itself a result.

\paragraph{Which observed trend to aim at.} Paired rates are $+0.458 \pm 0.048$ (1985--2006),
$+0.209 \pm 0.065$ (2006--2024) and $+0.312 \pm 0.035$ (1985--2024). The observed deceleration
factor of $2.19$ is steeper than pure approach-to-equilibrium predicts ($1.18$--$1.55$), which
would be the signature of 1985 reading low --- but the two rates differ by only $\approx 2\sigma$,
so the evidence is weak. We adopt $+0.312 \pm 0.035$ as the primary target (best determined,
spans the whole record) and report $+0.209$ as the trusted-campaign sensitivity.

% =====================================================================
\section{Verification before any recalibration}\label{sec:verify}

\begin{enumerate}
\item Transform round-trip: \texttt{to\_unconstrained(to\_original(x)) == x} for the new pair.
\item Jacobian checked numerically against a finite-difference determinant.
\item \texttt{preflight\_prior\_pushforward.R} at full $N$: forward blow-up rate should remain
$\approx 0\,\%$, as it did on the current target. This gate caught the \texttt{beta2} detonation
and is the natural one here.
\item Prior pushforward on $J_{1917}$ itself: confirm the implied 1917 flux distribution sits
inside the NPP envelope without the guard firing appreciably.
\item ESS on $R$ after the run --- a parameter against a boundary mixes worse, the caveat that
applied to $\sinput$ under \texttt{flux\_pair}.
\end{enumerate}

% =====================================================================
\section{Plan for the next calibration}

\subsection*{Established --- goes in regardless}
\begin{enumerate}
\item \textbf{Log-normal likelihood.} Confirmed on TP2: bias $+12.7 \rightarrow -0.9$
      tC\,ha$^{-1}$, median log-residual $-0.205 \rightarrow -0.023$. Implemented as
      \texttt{HIKET\_LOGNORMAL\_LIK}, default off. This is the single largest correction and it
      is not speculative.
\item \textbf{the common anchor --- common pre-run anchor.} Reference both ends to $\Jto$ so the sampled ratio
      \emph{is} the 1917/1985 ratio. A defect correction; neutral for fit.
\item \textbf{plain priors --- collapse the auxiliary priors} to two ordinary lognormals. Paired
      pre-flight shows identical blow-up rates with and without \texttt{flux\_pair} in all six
      models, and $0\,\%$ prior mass above the NPP ceiling.
\end{enumerate}

\subsection*{Aimed at the trend --- family-specific}
\begin{enumerate}\setcounter{enumi}{3}
\item \textbf{the ratio prior --- growing-stock prior on $R$}, for the \textbf{Yasso family}. Necessary and
      sufficient there: their $\tau$ is already right and only the starting position is wrong.
      Expected to convert the reported source back into a sink.
\item \textbf{the free slow rate --- relax the slow-rate anchor}, for \textbf{TP2/TP3}. The only lever on their
      trend; irrelevant to Yasso, whose rates are fixed. \emph{Under test.}
\end{enumerate}

\subsection*{Still open}
\begin{enumerate}\setcounter{enumi}{5}
\item \textbf{$\sinput$'s centre.} C4b set it at $1.30$ before the SOC target was corrected
      downward by a comparable factor. Deliberately not bundled: changing it alongside the
      error model would make the run uninterpretable. Note the log-normal refit left $\sinput$
      essentially unmoved, so this may matter less than it appeared.
\item \textbf{$\sobs$ fixed} (D4). A correctness issue independent of everything above, and
      partly superseded: under a log-normal the width no longer biases the location.
\item \textbf{Whether the NPP envelope should apply to the 1917 flux at all} --- a modern
      physiological ceiling extrapolated to a forest for which no NPP measurements exist.
\item \textbf{The observed target itself.} On the two LUKE-validated campaigns the observed
      trend is $+0.209$, not $+0.309$. Resolve before tuning models to close a gap that may be
      a third smaller than it looks.
\end{enumerate}

\subsection*{Verification gate before launch}
Transform round-trip; Jacobian against finite differences;
\texttt{preflight\_prior\_pushforward.R} at full $N$ ($\approx 0\,\%$ blow-ups expected); prior
pushforward on $J_{1917}$ itself; and ESS on $R$ after the run.

\subsection*{Pending at the time of writing}
Two short local runs were in progress: TP2 with log-normal $+$ widened slow-rate prior
(predicting $\alpha_H \rightarrow 0.012$--$0.025$ and the trend from $+0.116$ toward $+0.24$),
and Yasso20 with log-normal alone (testing whether the level fix generalises to a fixed-rate
model, and whether the inverted sink survives a correctly specified likelihood). Results to be
appended.

\end{document}
